% ----------------------
\subsection{Section 1 (1 November 1831)}
% ----------------------
	\label{DC1}
	
	% -----
	\subsubsection{Relation to section 67}
	% -----
	
		Section 1 was given at the same conference during which JS received \hyperref[DC67]{section 67}.
	
	% -----
	\subsubsection{Historical Background}
	% -----
		\label{DC1-HB}
		
		% --
		\paragraph{JSP D2}
		% --
		
			``At a conference of elders in Hiram, Ohio, on 1 November 1831, JS dictated a revelation designated as a preface for the Book of Commandments, a proposed compilation of JS's revelations. According to the minutes of the 1 November conference, this revelatory preface was `received by inspiration' during a recess between the morning and afternoon sessions. William E. McLellin, one of the conference participants, gave a more detailed account of the production of the preface fifty years later to William Kelley, an elder in the Reorganized Church of Jesus Christ of Latter Day Saints. According to Kelley's account of the conversation, McLellin said that he, Sidney Rigdon, and Oliver Cowdery had been given the assignment to write the preface to the Book of Commandments, but when they presented their draft to the conference, the `Conference picked it all to pieces' and requested that JS petition the Lord for a preface. After JS and the elders bowed in prayer, JS, who was `sitting by a window,' dictated the preface `by the Spirit,' while Rigdon served as scribe. `Joseph would deliver a few sentences and Sydney would write them down,' McLellin told Kelley, `then read them aloud, and if correct, then Joseph would proceed and deliver more.' In this way, `the preface was given.'%
				\footnote{%
					% \label{}%
					See below for the source for this account from McLellin and Kelley.
					}
			
			``As with prefaces to other published works, this revelation informed readers of the subject and purpose of the volume. Speaking in the voice of Deity, it told both the church and the world why God revealed these commandments to JS and painted an apocalyptic picture of the wrath that God would unleash against the wicked unless they responded to the new revelations by repenting. The revelations were thus a `voice of warning' to the world to prepare them for Jesus Christ's second coming.
			
			``The original manuscript of the preface is no longer extant. John Whitmer copied the revelation into Revelation Book 1 likely between 12 November 1831, the date of the revelation that precedes it in the revelation book, and 20 November, the date that he and Cowdery left for Missouri with the revelation book. The revelation was printed as the preface to the Book of Commandments about a year later'' (JSP D2:103--104).
			
		% --
		\paragraph{Far West Record}
		% --
		
			``Br Sidney Rigdon appointed moderator \& Oliver Cowdery Clerk br Oliver Cowdery made a request desiring the mind of the Lord through this conference of Elders to know how many copies of the Book of commandments it was the will of the Lord should be published in the first edition of that work. Voted that there be ten thousand \sout{struck} copies struck. ...
			
			``Preface received by inspiration...
			
			``Br. Joseph Smith jr. said that inasmuch as the Lord had bestowed a great blessing upon us in giving commandments and revelations, asked the conference what testimony they were willing to attach to these commandments which should shortly be sent to the world. A number of the brethren arose and said, that they were willing to testify to the world that they knew that they were of the Lord.
			
			``Revelation received relative to the same...
			
			``The Revelation of last evening read by the moderator the brethren then arose in turn and bore witness to the truth of the Book of Commandments. After which br. Joseph Smith jr. arose \& expressed his feelings \& gratitude concerning the commandment \& Preface recieved yesterday...(\href{https://www.josephsmithpapers.org/paper-summary/minute-book-2/17}{\textit{Far West Record}}, 15--16).
			
		% --
		\paragraph{Letter From Elder W. H. Kelley (Saints' Herald)}
		% --
		
			``In August last I left home to attend the Semi-Annual Conference, appointed to convene near Council Bluffs, Iowa...In Chicago I was joined by our excellent brother, George A. Blakeslee...

			``In company with Bro. Warnky and Dr. Wm. E. McLellin, we visited the 1temple lot' lying west of the [Independence] Court House. The doctor was able to point of the identical spot where Joseph stood when he first visited it, and which is the place of the corner stone....I gleaned from him and the records in his possession, the following items: He was present when the preface to the Book of Commandments was given, and says that Sidney Rigdon wrote it down as it was doctated by Joseph. A committee had been appointed to draft a preface, consisting of himself, O. Cowdery and, I think, Sidney Rigdon, but when they made their report, the `Conference picked it all to pieces.' The Conference then requested Joseph to enquire of the Lord about it, and he said that he would if the people would bow in prayer with him. This they did and Joseph prayed.

			``When they arose, Joseph dictated by the Spirit the preface found in the Book of Doctrine and Covenants while sitting by a window of the room in which the conference was sitting; and Sidney Rigdon wrote it down. Joseph would deliver a few sentences and Sidney would write them down, then read them aloud, and if correct, then Joseph would proceed and deliver more, and by this process the preface was given.

			``In reply to the question, `did Joseph seem to be inspired at the time? that is, did any thing of unusual character appear to be moving him?' he said, `There was something a hold of him'...'' (\textit{Saints' Herald}, 1 March 1882).%
				\footnote{%
					% \label{}%
					I copied and pasted this text from \href{http://www.sidneyrigdon.com/dbroadhu/IA/sain1882.htm}{this link}; I have not been able to find an original copy of the publication to check it against.
					}
					
		% --
		\paragraph{HDDC}
		% --
		
			See HDDC 21--26 (PDF 51--56) for more discussion of the conferences during which JS received section 1.
		
	% -----
	\subsubsection{Versions}
	% -----
		\label{DC1-V}
		
		See the \href{https://drive.google.com/file/d/1goAvNz8jS4R8slYfMG_db55Ty2z_vmgK/view?usp=sharing}{sections comparison} document.
	
		\begin{compactdesc}
			\item[JSP] \hyperref[EMS-1831-JS-1-NOV-1831-B]{1 November 1831--B}
			\item[EVMS] 1:10 (March 1833), 78
			\item[BC] Chapter 1, 3--6
			\item[DC 1835] Section 1, 75--77
			\item[EVMSR] 1:10 (``March 1833''; May 1836), 155--156
			\item[TS] 3:5 (1 January 1842), 639--640
			\item[TS] 5:7 (1 April 1844), 483--484%
				\footnote{%
					% \label{}%
					The second printing of section 1 in the TS was as part of JS's history, which was being printed in installments.
					}
			\item[DC 1844] Section 1, 87--91
			\item[MS] 5:11 (April 1845), 164--165%
				\footnote{%
					% \label{}%
					This printing of section 1 was also as part of JS's history.
					}
		\end{compactdesc}

	% -----
	\subsubsection{Section 1:1} % *DC1-1 
	% -----
		\label{DC1-1}

		\hyperref[HEARKEN]{HEARKEN}, O ye people of my church, saith the voice of him who dwells on high, and whose eyes are upon all men; yea, verily I say: \hyperref[HEARKEN]{Hearken} ye people from afar; and ye that are upon the islands of the sea, listen together.

		% \begin{framed}
		% \scriptsize
			% \begin{compactdesc}
				% \item[JSP] 
				% % \item[BC] 
				% % \item[]
			% \end{compactdesc}
		% \end{framed}

	% -----
	\subsubsection{Section 1:2} % *DC1-2 
	% -----
		\label{DC1-2}

		For verily the voice of the Lord is unto all men, and there is none to escape; and there is no eye that shall not see, neither ear that shall not hear, neither heart that shall not be penetrated.

		% \begin{framed}
		% \scriptsize
			% \begin{compactdesc}
				% \item[JSP] 
				% % \item[BC] 
				% % \item[]
			% \end{compactdesc}
		% \end{framed}

	% -----
	\subsubsection{Section 1:3} % *DC1-3 
	% -----
		\label{DC1-3}

		And the rebellious shall be pierced with much sorrow;%
			\footnote{%
				% \label{}%
				See \hyperref[1TIMOTHY-6-10]{1 Timothy 6:10}.
				}
		for their iniquities shall be spoken upon the housetops, and their secret acts shall be revealed.

		% \begin{framed}
		% \scriptsize
			% \begin{compactdesc}
				% \item[JSP] 
				% % \item[BC] 
				% % \item[]
			% \end{compactdesc}
		% \end{framed}

	% -----
	\subsubsection{Section 1:4} % *DC1-4 
	% -----
		\label{DC1-4}

		And the voice of warning shall be unto all people, by the mouths of my disciples, whom I have chosen in these last days.

		% \begin{framed}
		% \scriptsize
			% \begin{compactdesc}
				% \item[JSP] 
				% % \item[BC] 
				% % \item[]
			% \end{compactdesc}
		% \end{framed}

	% -----
	\subsubsection{Section 1:5} % *DC1-5 
	% -----
		\label{DC1-5}

		And they shall go forth and none shall stay them, for I the Lord have commanded them.

		% \begin{framed}
		% \scriptsize
			% \begin{compactdesc}
				% \item[JSP] 
				% % \item[BC] 
				% % \item[]
			% \end{compactdesc}
		% \end{framed}

	% -----
	\subsubsection{Section 1:6} % *DC1-6 
	% -----
		\label{DC1-6}

		Behold, this is mine authority, and the authority of my servants, and my preface unto the book of my commandments, which I have given them to publish unto you, O inhabitants of the earth.

		% \begin{framed}
		% \scriptsize
			% \begin{compactdesc}
				% \item[JSP] 
				% % \item[BC] 
				% % \item[]
			% \end{compactdesc}
		% \end{framed}

	% -----
	\subsubsection{Section 1:7} % *DC1-7 
	% -----
		\label{DC1-7}

		Wherefore, fear and tremble,%
			\footnote{%
				% \label{}%
				The ``fear and tremble'' combination is common both in the OT/NT and in the BOM.
				}
		O ye people, for what I the Lord have decreed in them shall be fulfilled.

		% \begin{framed}
		% \scriptsize
			% \begin{compactdesc}
				% \item[JSP] 
				% % \item[BC] 
				% % \item[]
			% \end{compactdesc}
		% \end{framed}

	% -----
	\subsubsection{Section 1:8} % *DC1-8 
	% -----
		\label{DC1-8}

		And verily I say unto you, that they who go forth, bearing these tidings unto the inhabitants of the earth, to them is power given to \hyperref[SEAL]{seal} both on earth and in heaven, the unbelieving and rebellious;

		% \begin{framed}
		% \scriptsize
			% \begin{compactdesc}
				% \item[JSP] 
				% % \item[BC] 
				% % \item[]
			% \end{compactdesc}
		% \end{framed}

	% -----
	\subsubsection{Section 1:9} % *DC1-9 
	% -----
		\label{DC1-9}

		Yea, verily, to seal them up unto the day when the wrath of God shall be poured out upon the wicked without measure---

		% \begin{framed}
		% \scriptsize
			% \begin{compactdesc}
				% \item[JSP] 
				% % \item[BC] 
				% % \item[]
			% \end{compactdesc}
		% \end{framed}

	% -----
	\subsubsection{Section 1:10} % *DC1-10 
	% -----
		\label{DC1-10}

		Unto the day when the Lord shall come to recompense unto every man according to his work, and measure to every man according to the measure which he has measured to his fellow man.%
			\footnote{%
				% \label{}%
				We ourselves become the standard against which we are judged, in a way. It reminds me of Plato, quoting Protagoras, in the \textit{Theaetetus}: ``Man is the measure of all things: of the things which are, that they are, and of the things which are not, that they are not'' (152a). The most important scriptural connection here is \hyperref[MATTHEW-7-2]{Matthew 7:2} (see also \hyperref[MARK-4-24]{Mark 4:24} and \hyperref[LUKE-6-38]{Luke 6:38}).
				}

		% \begin{framed}
		% \scriptsize
			% \begin{compactdesc}
				% \item[JSP] 
				% % \item[BC] 
				% % \item[]
			% \end{compactdesc}
		% \end{framed}

	% -----
	\subsubsection{Section 1:11} % *DC1-11 
	% -----
		\label{DC1-11}

		Wherefore the voice of the Lord is unto the ends of the earth, that all that will hear may hear:

		% \begin{framed}
		% \scriptsize
			% \begin{compactdesc}
				% \item[JSP] 
				% % \item[BC] 
				% % \item[]
			% \end{compactdesc}
		% \end{framed}

	% -----
	\subsubsection{Section 1:12} % *DC1-12 
	% -----
		\label{DC1-12}

		Prepare ye, prepare ye for that which is to come, for the Lord is nigh;

		% \begin{framed}
		% \scriptsize
			% \begin{compactdesc}
				% \item[JSP] 
				% % \item[BC] 
				% % \item[]
			% \end{compactdesc}
		% \end{framed}

	% -----
	\subsubsection{Section 1:13} % *DC1-13 
	% -----
		\label{DC1-13}

		And the anger of the Lord is kindled, and his sword is bathed in heaven, and it shall fall upon the inhabitants of the earth.%
			\footnote{%
				% \label{}%
				The ``bathed in heaven'' part, and the bit about the sword falling, is a reference to \hyperref[ISAIAH-34-5]{Isaiah 34:5}.
				}

		% \begin{framed}
		% \scriptsize
			% \begin{compactdesc}
				% \item[JSP] 
				% % \item[BC] 
				% % \item[]
			% \end{compactdesc}
		% \end{framed}

	% -----
	\subsubsection{Section 1:14} % *DC1-14 
	% -----
		\label{DC1-14}

		And the arm of the Lord shall be revealed; and the day cometh that they who will not hear the voice of the Lord, neither the voice of his servants, neither give heed to the words of the prophets and apostles, shall be cut off from among the people;

		% \begin{framed}
		% \scriptsize
			% \begin{compactdesc}
				% \item[JSP] 
				% % \item[BC] 
				% % \item[]
			% \end{compactdesc}
		% \end{framed}

	% -----
	\subsubsection{Section 1:15} % *DC1-15 
	% -----
		\label{DC1-15}

		For they have strayed from mine ordinances, and have broken mine everlasting covenant;%
		\footnote{%
			% \label{}%
			The language here seems to come from \hyperref[ISAIAH-24-5]{Isaiah 24:5}, where people have ``transgressed the laws'' and ``changed the ordinance'' and ``broken the everlasting covenant.'' ``Everlasting covenant'' appears 35 times in the scriptures (see occurrences \hyperref[COVENANT-PHRASES]{here}); 14 in the OT, 1 in the NT (\hyperref[HEBREWS-13-20]{Hebrews 13:20}), and 20 in the DC. It appears in significant places in the OT as well, including when God makes a covenant with Noah, when God makes a covenant with Abraham, and so on.
			} 

		% \begin{framed}
		% \scriptsize
			% \begin{compactdesc}
				% \item[JSP] 
				% % \item[BC] 
				% % \item[]
			% \end{compactdesc}
		% \end{framed}

	% -----
	\subsubsection{Section 1:16} % *DC1-16 
	% -----
		\label{DC1-16}

		They seek not the Lord to establish his righteousness, but every man walketh in his own way, and after the image of his own god, whose image is in the likeness of the world, and whose substance is that of an idol, which waxeth old and shall perish in Babylon, even Babylon the great, which shall fall.

		% \begin{framed}
		% \scriptsize
			% \begin{compactdesc}
				% \item[JSP] 
				% % \item[BC] 
				% % \item[]
			% \end{compactdesc}
		% \end{framed}

	% -----
	\subsubsection{Section 1:17} % *DC1-17 
	% -----
		\label{DC1-17}

		Wherefore, I the Lord, knowing the calamity which should come upon the inhabitants of the earth, called upon my servant Joseph Smith, Jun., and spake unto him from heaven, and gave him commandments;

		% \begin{framed}
		% \scriptsize
			% \begin{compactdesc}
				% \item[JSP] 
				% % \item[BC] 
				% % \item[]
			% \end{compactdesc}
		% \end{framed}

	% -----
	\subsubsection{Section 1:18} % *DC1-18 
	% -----
		\label{DC1-18}

		And also gave commandments to others, that they should proclaim these things unto the world; and all this that it might be fulfilled, which was written by the prophets---

		% \begin{framed}
		% \scriptsize
			% \begin{compactdesc}
				% \item[JSP] 
				% % \item[BC] 
				% % \item[]
			% \end{compactdesc}
		% \end{framed}

	% -----
	\subsubsection{Section 1:19} % *DC1-19 
	% -----
		\label{DC1-19}

		The weak things of the world shall come forth and break down the mighty and strong ones, that man should not counsel his fellow man, neither trust in the arm of flesh---

		% \begin{framed}
		% \scriptsize
			% \begin{compactdesc}
				% \item[JSP] 
				% % \item[BC] 
				% % \item[]
			% \end{compactdesc}
		% \end{framed}

	% -----
	\subsubsection{Section 1:20} % *DC1-20 
	% -----
		\label{DC1-20}

		But that every man might speak in the name of God the Lord, even the Savior of the world;

		% \begin{framed}
		% \scriptsize
			% \begin{compactdesc}
				% \item[JSP] 
				% % \item[BC] 
				% % \item[]
			% \end{compactdesc}
		% \end{framed}

	% -----
	\subsubsection{Section 1:21} % *DC1-21 
	% -----
		\label{DC1-21}

		That faith also might increase in the earth;

		% \begin{framed}
		% \scriptsize
			% \begin{compactdesc}
				% \item[JSP] 
				% % \item[BC] 
				% % \item[]
			% \end{compactdesc}
		% \end{framed}

	% -----
	\subsubsection{Section 1:22} % *DC1-22 
	% -----
		\label{DC1-22}

		That mine everlasting covenant might be established;

		% \begin{framed}
		% \scriptsize
			% \begin{compactdesc}
				% \item[JSP] 
				% % \item[BC] 
				% % \item[]
			% \end{compactdesc}
		% \end{framed}

	% -----
	\subsubsection{Section 1:23} % *DC1-23 
	% -----
		\label{DC1-23}

		That the fulness of my gospel might be proclaimed by the weak and the simple unto the ends of the world, and before kings and rulers.

		% \begin{framed}
		% \scriptsize
			% \begin{compactdesc}
				% \item[JSP] 
				% % \item[BC] 
				% % \item[]
			% \end{compactdesc}
		% \end{framed}

	% -----
	\subsubsection{Section 1:24} % *DC1-24 
	% -----
		\label{DC1-24}

		Behold, I am God and have spoken it; these commandments are of me, and were given unto my servants in their weakness,%
			\footnote{%
				% \label{}%
				The BOM has many uses of the word ``weakness'' which are relevant here, I think; Nephi uses the term several times in talking about what he is writing, for example. We can make the connection to JS---as a servant he had this ``weakness,'' which we should expect to result in an imperfect process of communication with God.
				}
		after the manner of their language,%
			\footnote{%
				% \label{}%
				This part of the verse is very close to \hyperref[2NEPHI-31-3]{2 Nephi 31:3}.
				}
		that they might come to understanding.%
			\footnote{%
				% \label{}%
				Actually reaching a level of understanding is key; see \hyperref[DC68-25]{DC 68:25} for example.
				}

		\begin{framed}
		\scriptsize
			\begin{compactdesc}
				\item[JSP] Behold I am God \& have spoken it these \sout{are} commandments are of me \& were given unto my Servents in their weakness after the manner of their Language that they might come to understanding
				\item[EMS] Behold I am God and have spoken it: These commandments are of me, and were given unto my servants in their weakness, after the manner of their language, that they might come to understanding;
				\item[BC] Behold I am God and have spoken it: these commandments are of me, and were given unto my servants in their weakness, after the manner of their language, that they might come to understanding;
				\item[DC 1835] Behold I am God and have spoken it: these commandments are of me, and were given unto my servants in their weakness, after the manner of their language, that they might come to understanding;
				\item[DC 1844] Behold I am God and have spoken it: these commandments are of me, and were given unto my servants in their weakness, after the manner of their language, that they might come to understanding;
			\end{compactdesc}
		\end{framed}

	% -----
	\subsubsection{Section 1:25} % *DC1-25 
	% -----
		\label{DC1-25}

		And inasmuch as they erred it might be made known;

		% \begin{framed}
		% \scriptsize
			% \begin{compactdesc}
				% \item[JSP] 
				% % \item[BC] 
				% % \item[]
			% \end{compactdesc}
		% \end{framed}

	% -----
	\subsubsection{Section 1:26} % *DC1-26 
	% -----
		\label{DC1-26}

		And inasmuch as they sought wisdom they might be instructed;

		% \begin{framed}
		% \scriptsize
			% \begin{compactdesc}
				% \item[JSP] 
				% % \item[BC] 
				% % \item[]
			% \end{compactdesc}
		% \end{framed}

	% -----
	\subsubsection{Section 1:27} % *DC1-27 
	% -----
		\label{DC1-27}

		And inasmuch as they sinned they might be chastened, that they might repent;

		% \begin{framed}
		% \scriptsize
			% \begin{compactdesc}
				% \item[JSP] 
				% % \item[BC] 
				% % \item[]
			% \end{compactdesc}
		% \end{framed}

	% -----
	\subsubsection{Section 1:28} % *DC1-28 
	% -----
		\label{DC1-28}

		And inasmuch as they were humble they might be made strong, and blessed from on high, and receive knowledge from time to time.

		% \begin{framed}
		% \scriptsize
			% \begin{compactdesc}
				% \item[JSP] 
				% % \item[BC] 
				% % \item[]
			% \end{compactdesc}
		% \end{framed}

	% -----
	\subsubsection{Section 1:29} % *DC1-29 
	% -----
		\label{DC1-29}

		And after having received the record of the Nephites, yea, even my servant Joseph Smith, Jun., might have power to translate through the mercy of God, by the power of God, the Book of Mormon.

		% \begin{framed}
		% \scriptsize
			% \begin{compactdesc}
				% \item[JSP] 
				% % \item[BC] 
				% % \item[]
			% \end{compactdesc}
		% \end{framed}

	% -----
	\subsubsection{Section 1:30} % *DC1-30 
	% -----
		\label{DC1-30}

		And also those to whom these commandments were given, might have power to lay the foundation of this church, and to bring it forth out of obscurity and out of darkness, the only true and living church upon the face of the whole earth, with which I, the Lord, am well pleased, speaking unto the church collectively and not individually---

		% \begin{framed}
		% \scriptsize
			% \begin{compactdesc}
				% \item[JSP] 
				% % \item[BC] 
				% % \item[]
			% \end{compactdesc}
		% \end{framed}

	% -----
	\subsubsection{Section 1:31} % *DC1-31 
	% -----
		\label{DC1-31}

		For I the Lord cannot look upon sin with the least degree of allowance;%
			\footnote{%
				% \label{}%
				See \hyperref[ALMA-45-16]{Alma 45:16}.
				}

		% \begin{framed}
		% \scriptsize
			% \begin{compactdesc}
				% \item[JSP] 
				% % \item[BC] 
				% % \item[]
			% \end{compactdesc}
		% \end{framed}

	% -----
	\subsubsection{Section 1:32} % *DC1-32 
	% -----
		\label{DC1-32}

		Nevertheless, he that repents and does the commandments of the Lord shall be forgiven;

		% \begin{framed}
		% \scriptsize
			% \begin{compactdesc}
				% \item[JSP] 
				% % \item[BC] 
				% % \item[]
			% \end{compactdesc}
		% \end{framed}

	% -----
	\subsubsection{Section 1:33} % *DC1-33 
	% -----
		\label{DC1-33}

		And he that repents not, from him shall be taken even the light which he has received;%
			\footnote{%
				% \label{}%
				In my mind the central place for this idea is \hyperref[2NEPHI-28-30]{2 Nephi 28:30}, though the principle there is a little different (or maybe more general---there are many ways of saying ``we've had enough,'' and unrighteousness is just one of them) from the one here.
				}
		for my Spirit shall not always strive with man, saith the Lord of Hosts.

		% \begin{framed}
		% \scriptsize
			% \begin{compactdesc}
				% \item[JSP] 
				% % \item[BC] 
				% % \item[]
			% \end{compactdesc}
		% \end{framed}

	% -----
	\subsubsection{Section 1:34} % *DC1-34 
	% -----
		\label{DC1-34}

		And again, verily I say unto you, O inhabitants of the earth: I the Lord am willing to make these things known unto all flesh;

		% \begin{framed}
		% \scriptsize
			% \begin{compactdesc}
				% \item[JSP] 
				% % \item[BC] 
				% % \item[]
			% \end{compactdesc}
		% \end{framed}

	% -----
	\subsubsection{Section 1:35} % *DC1-35 
	% -----
		\label{DC1-35}

		For I am no respecter of persons,%
		\footnote{%
			% \label{}%
			This is a construction that actually appears in both the OT and NT; it's a very ancient idea. For example, see \hyperref[DEUTERONOMY-1-17]{Deuteronomy 1:17} and \hyperref[DEUTERONOMY-16-19]{16:19}, \hyperref[2SAMUEL-14-14-]{2 Samuel 14:14} (this one specifically mentions God), \hyperref[ACTS-10-34]{Acts 10:34}, \hyperref[ROMANS-2-11]{Romans 2:11}, \hyperref[1PETER-1-17]{1 Peter 1:17}, and so on.
			} 
		and will that all men shall know that the day speedily cometh; the hour is not yet, but is nigh at hand, when peace shall be taken from the earth, and the devil shall have power over his own dominion.

		% \begin{framed}
		% \scriptsize
			% \begin{compactdesc}
				% \item[JSP] 
				% % \item[BC] 
				% % \item[]
			% \end{compactdesc}
		% \end{framed}

	% -----
	\subsubsection{Section 1:36} % *DC1-36 
	% -----
		\label{DC1-36}

		And also the Lord shall have power over his saints, and shall reign in their midst, and shall come down in judgment upon Idumea, or the world.

		% \begin{framed}
		% \scriptsize
			% \begin{compactdesc}
				% \item[JSP] 
				% % \item[BC] 
				% % \item[]
			% \end{compactdesc}
		% \end{framed}

	% -----
	\subsubsection{Section 1:37} % *DC1-37 
	% -----
		\label{DC1-37}

		Search these commandments,%
			\footnote{%
				% \label{}%
				I read this injunction as being like \hyperref[JOHN-5-39]{John 5:39}, where the imperative verb is a form of \hyperref[ERAUNAO]{\grl{ἐραυνάω}}. Not just look through them---\textit{examine}, \textit{investigate}, \textit{seek knowledge of}.
				}
		for they are true and faithful,%
			\footnote{%
				% \label{}%
				See \hyperref[DC68-34]{DC 68:34}, given at the same conference as this section.
				}
		and the prophecies and promises which are in them shall all be fulfilled.

		% \begin{framed}
		% \scriptsize
			% \begin{compactdesc}
				% \item[JSP] 
				% % \item[BC] 
				% % \item[]
			% \end{compactdesc}
		% \end{framed}

	% -----
	\subsubsection{Section 1:38} % *DC1-38 
	% -----
		\label{DC1-38}

		What I the Lord have spoken, I have spoken, and I excuse not myself;%
			\footnote{%
				% \label{}%
				Check out meaning \#6 for ``excuse'' from the 1828 dictionary: ``to throw off an imputation by apology.'' 
				}
		and though the heavens and the earth pass away, my word shall not pass away, but shall all be fulfilled,%
			\footnote{%
				% \label{}%
				See \hyperref[MATTHEW-24-35]{Matthew 24:35} and notes there.
				}
		whether by mine own voice or by the voice of my servants, it is the same.%
			\footnote{%
				% \label{}%
				I think I used to take this verse as saying that it didn't matter whether it was the servants of God speaking, or God himself, it's the same thing. I don't think it actually says that, though. At least, the ``whether by mine own voice'' part at the end seems to fall just within the scope of the fulfillment of God's word spoken about immediately before that. It doesn't seem like it's supposed to have a more general scope than that.
				}

		\begin{framed}
		\scriptsize
			\begin{compactdesc}
				\item[JSP] what I the Lord have spoken I have spoken \& I excuse not myself \& though the Heaven \& <the> Earth pass away my word shall not pass away but shall all be fulfilled whether by mine own voice or by the voice of my Servants it is the same
				\item[BC] What I the Lord have spoken, I have spoken, and I excuse not myself, and though the heavens and the earth pass away, my word shall not pass away, but shall all be fulfilled, whether by mine own voice, or by the voice of my servants, it is the same:
				\item[DC 1835] What I the Lord have spoken, I have spoken, and I excuse not myself, and though the heavens and the earth pass away, my word shall not pass away, but shall all be fulfilled, whether by mine own voice, or by the voice of my servants, it is the same:
				\item[DC 1844] What I the Lord have spoken, I have spoken, and I excuse not myself, and though the heavens and the earth pass away, my word shall not pass away, but shall all be fulfilled, whether by mine own voice, or by the voice of my servants, it is the same:
			\end{compactdesc}
		\end{framed}

	% -----
	\subsubsection{Section 1:39} % *DC1-39 
	% -----
		\label{DC1-39}

		For behold, and lo, the Lord is God, and the Spirit beareth record, and the record is true, and the truth abideth forever and ever. Amen.

		% \begin{framed}
		% \scriptsize
			% \begin{compactdesc}
				% \item[JSP] 
				% % \item[BC] 
				% % \item[]
			% \end{compactdesc}
		% \end{framed}